\section{Multivariable Calculus}
\subsection{Operators}
\begin{itemize}
\item Dot product (點積) operator: $\cdot$
\item Cross product (叉積) operator: $\times$
\item Gradient (梯度) operator: $\nabla$ or $\operatorname{grad}$
\item Jacobian matrix (雅可比矩陣): $\nabla$, or $\mathbf{J}(\mathbf{F})$, $J(\mathbf{F})$, $\mb{J}_{\mb{F}}$, or $J_{\mathbf{F}}$ for operand $\mb{F}$.
\item Gram matrix, Gramian matrix, or Gramian: $G(\mb{F})$, $G_{\mb{F}}$ for operand $\mb{F}$.
\item Divergence (散度) operator: $\nabla \cdot$ or $\operatorname{div}$
\item Curl (旋度) operator: $\nabla \times$ or $\operatorname{curl}$
\item Directional derivative (方向導數) operator: $(\mb{F}\cdot\nabla)\mb{A}$ or $D_{\mb{F}}\mb{A}$ for directional derivative of scalar field $\mb{A}$ along a vector $\mb{F}$, typically unit vector and $D_{\mb{F}}\mb{A}$ is typically used, or vector field $\mb{F}$ and $(\mb{F}\cdot\nabla)\mb{A}$ is typically used.
\item Hessian (黑塞) (matrix): $D^2f$, $\mb{H}(f)$, $H(f)$, $\mb{H}_f$, $H_f$, or $\nabla\otimes\nabla f$ for operand $f$.
\item Laplace operator (拉普拉斯算子): $\nabla\cdot\nabla$, $\nabla^2$, or $\Delta$. For clarity, we use $\nabla\cdot\nabla$ here.
\item Line/path integral operator: $\int$
\item Double/surface integral operator: $\iint$
\item Triple/volume integration operator: $\iiint$
\item Closed line integral operator: $\oint$
\item Closed surface/double integral operator: $\oiint$
\item $\int\mathbf{F}\cdot\mathrm{d}\mathbf{S}$ is used as a shorthand for $\int(\mathbf{F}\cdot\mathbf{\hat{n}})\,\mathrm{d}S$, where $\hat{n}$ is the outward pointing unit normal at almost each point on $S$.
\end{itemize}
\subsection{Convention}
If not otherwise specified:
\begin{itemize}
\item The domain of the funcitons or maps below are subsets of a Euclidean vector space. If not otherwise specified, the coordinates are the Cartesian coordinates, the norms are the Euclidean norms, and the measures are the Lebesgue measures.
\item $\mathbf{0}$ or $0$ refers to the zero tensor (零張量) in the interested Euclidean tensor space $V$, that is, it satisfies 
\[\forall\mathbf{v}\in V:\,\mathbf{v}+\mathbf{0}=\mathbf{v}.\]
\item Unit vector (單位向量): $\mathbf{e}_i$ is the unit vector in the $i$th direction, i.e., a vector with zero norm.
\item Independent variable vector: $\mathbf{x}=(x_1,x_2,\dots,x_n)$
\item Vector fields: $\mathbf{F}(\mathbf{x}) = \sum_{i=1}^n F_i(\mathbf{x}) \mathbf{e}_i$、$\mathbf{G}$
\item Scalar fields: $A(\mathbf{x})$、$B(\mathbf{x})$
\item Tensor fields: $f(\mathbf{x})$、$g(\mathbf{x})$
\item Three-dimensional tensor space field: $\mathbf{T}(\mathbf{x})$
\item The $i$th component of the map $f$: $f_i$
\end{itemize}
\subsection{First order}
\subsubsection{Gradient (梯度) or first derivative, and Jacobian matrix}
\[
\nabla f = \begin{pmatrix}\qty(\pdv{f}{x_1})^T & \qty(\pdv{f}{x_2})^T & \dots  & \qty(\pdv{f}{x_n})^T\end{pmatrix}
\]
The gradient of a scalar field is a vector field, the gradient of a vector field is a second-order tensor (matrix) field, and the gradient of a $k$-order tensor field is a $k+1$-order tensor field. 

In particular, the gradient of a scalar field $A$ is
\[
\nabla A = \sum_{i=1}^n \pdv{A}{x_i}e_i.
\]
And the gradient of a vector field $\mathbf{F}$, also called the Jacobian matrix (雅可比矩陣) of $\mb{F}$, is a matrix such that the $( i,j )$th entry is
\[(J(\mb{F}))_{ij}=\frac{\partial F_i}{\partial x_j}.\]
The determinant of Jacobian matrix is called Jacobian determinant, or Jacobian for short.
\sssc{Gram matrix, Gramian matrix, or Gramian}
The Gram matrix, Gramian matrix, or Gramian of $G$ of $\mb{F}$ is defined as an $n\times n$ matrix where
\[G_{ij}=\pdv{\mb{F}}{x_i}\cdot\pdv{\mb{F}}{x_j},\]
that is,
\[G=J^{\top}J,\]
where $J$ is the Jacobian matrix of $\mb{F}$.
 
The Gram determinant or Gramian is the determinant of the Gram matrix.
\subsubsection{Divergence (散度)}
\[
\nabla \cdot f = \sum_{i=1}^n\frac{\partial f_i}{\partial x_i}
\]
The divergence of a vector field is a scalar field, the gradient of a second-order tensor (i.e. matrix) field is a vector field, and the divergence of a $k+1$-order tensor field is a $k$-order tensor field.
\subsubsection{Curl (旋度)}
The curl is only defined on three-dimensional vector field.
\[
\nabla \times \mathbf{T} = 
\begin{pmatrix}
\mathbf{e}_1 & \mathbf{e}_2 & \mathbf{e}_3 \\
\frac{\partial}{\partial x_1} & \frac{\partial}{\partial x_2} & \frac{\partial}{\partial x_3} \\
T_1 & T_2 & T_3 \\
\end{pmatrix}
\]
The curl of a three-dimensional vector field is a three-dimensional vector field.
\subsubsection{Directional derivative (方向導數)}
For $\mb{A}\colon D\subseteq\bbR^n\to\bbR$ and $\mb{F}\colon E\subseteq\bbR^n\to\bbR^n$,
\[(\mb{F}\cdot\nabla)\mb{A}=D_{\mb{F}}\mb{A}\coloneq\lim_{h\to 0}\frac{\mb{A}(\mb{x}+h\mb{F}(\mb{x}))-\mb{A}(\mb{x})}{h}.\]
For points where $\mb{A}$ is differentiable,
\[(\mb{F}\cdot\nabla)\mb{A}=\sum_{i=1}^nF_i\pdv{\mb{A}}{x_i}=\mb{F}\cdot\nabla\mb{A}.\]
In $D_{\mb{F}}\mb{A}$ notation, $\mb{F}$ is often constant and often unit vector.
\subsection{Higher order}
\subsubsection{Hessian (黑塞) (matrix) or second derivative}
The Hessian or second derivative $H_f$ of field $f$ is defined as a matrix such that
\[(H(f))_{ij}=\frac{\partial f}{\partial x_i\partial x_j},\]
that is,
\[H(f)=\nabla\qty((\nabla f)^\top).\]
\subsubsection{Laplace operator (拉普拉斯算子)}
\[\nabla\cdot\nabla f = \nabla \cdot (\nabla f) = \tr(H(f))=\sum_{i=1}^n\frac{\partial^2 f}{\partial x_i^{\phantom{i}2}}\]
The Laplace operator applied to a tensor field is a tensor field of the same order and same dimension.
\subsubsection{Poisson's equation (卜瓦松 or 帕松 or 泊松方程)}
\[
\nabla\cdot\nabla A = B(\mathbf{x})
\]
\subsubsection{Laplace's equation (拉普拉斯方程)}
\[
\nabla\cdot\nabla A = 0
\]
A real function $A$ with real independent variables that is second-order differentiable for all independent variables is called a harmonic function if $A$ satisfies Laplace's equation.
\subsubsection{Multi-index notation (多重指標記號)}
Suppose there are \( n \) variables \( x_1, x_2, \dots, x_n \), then a multi-index $\alpha$ is a vector of \( n \) nonnegative integers: 
\[
\alpha = (\alpha_1, \alpha_2,\ldots,\alpha_n), \quad \text{where\ } \alpha_i \in \mathbb{N}_0.
\]
Define: 
\begin{itemize}
\item Norm or length $|\alpha|$ or \( \|\alpha\| \): 
\[
|\alpha| = \alpha_1 + \alpha_2 + \ldots\alpha_n.
\]
\item Factorial \( \alpha! \): 
\[
\alpha! = \alpha_1! \cdot \alpha_2! \cdot \ldots \alpha_n!.
\]
\item Power \( \mathbf{x}^\alpha \): 
If \( \mathbf{x} = (x_1, x_2,\ldots, x_n) \), then
\[
\mathbf{x}^\alpha = x_1^{\alpha_1} \cdot x_2^{\alpha_2} \ldots  x_n^{\alpha_n}.
\]
\item Order:
\[(\beta_1, \beta_2,\ldots,\beta_n)\leq(\alpha_1, \alpha_2,\ldots, \alpha_n)\coloneq\beta_1\leq\alpha_1\land\beta_2\leq\alpha_2\land\ldots\beta_n\leq\alpha_n\]
\[(\beta_1, \beta_2,\ldots,\beta_n)<(\alpha_1, \alpha_2,\ldots, \alpha_n)\coloneq\beta_1<\alpha_1\land\beta_2<\alpha_2\land\ldots\beta_n<\alpha_n\]
\item Higher-order mixed partial derivatives $D^\alpha f$ or $\partial^\alpha f$: 
\[
D^\alpha f = \frac{\partial^{|\alpha|} f}{\partial x_1^{\alpha_1} \partial x_2^{\alpha_2} \ldots  \partial x_n^{\alpha_n}}.
\]
\end{itemize}
\subsubsection{Higher-order derivative}
The $k$th order derivative of $\mathbf{F}:\,\mathbb{R}^n\to\mathbb{R}^m$, denoted as $\mathbf{F}^{(k)}(\mathbf{x})$ or $D^{k}\mathbf{F}(\mathbf{x})$, is a $(\mathbb{R}^n)^k\to\mathbb{R}$ function, where $(\mathbb{R}^n)^k$ is a Cartesian product of $k$ copies of $\mathbb{R}^n$ vector, that is,
\[D^{k}\mathbf{F}(\mathbf{x})=\sum_{|\alpha|=k} \left(D^\alpha \mathbf{F}(\mathbf{a})\right).\]
In particular, the first-order derivative of $\mathbf{F}$ is the gradient of it.
\subsubsection{Taylor series or Taylor expansion}
Let $f\colon D\subseteq\mathbb{R}^m\to\mathbb{R}$ be a function and $I\subseteq D$ be an open ball centered at $a$ such that $f\in C^{\infty}(I)$, then the Taylor series of $f$ at $a$ (or about $a$, near $a$, or centered at $a$) is
\[\sum_{\alpha\in\mathbb{N}_0^{\pht{0}m}}\frac{D^{\alpha}f(a)}{\alpha !}(x-a)^{\alpha},\]
of which convergence is not guaranteed, and convergence to $f$ is not guaranteed even if it's convergent.

Taylor series of at $0$ is called Maclaurin series 
or Maclaurin expansion.
\subsubsection{Taylor polynomial}
Let $f\colon D\subseteq\mathbb{R}^m\to\mathbb{R}$ be a function and $I\subseteq D$ be an open ball centered at $a$ where $f$ is $k$ times differentiable, then the $k$-th degree (or order) Taylor polynomial (or approximation) of $f$ at $a$ (or about $a$, near $a$, or centered at $a$) is
\[T_n(x)=\sum_{\alpha\in\mathbb{N}_0^{\pht{0}m}\land|\alpha|\leq k}\frac{D^{\alpha}f(a)}{\alpha !}(x-a)^{\alpha},\]
in which the first degree one is also called linear approximation, or linearization, or in three dimensions, tangent plane (approximation), and
\[R_n(x)=f(x)-T_n(x)\]
is called remainder.

Taylor polynomial of at $0$ is called Maclaurin polynomial.
\ssc{Graph}
\sssc{Level curve, level slice, or contour line}
For a real-valued function $f$ of two variables, the level curve of $f$ at level $k\in\bbR$ is defined as the set of $\mb{x}$ such that
\[f(\mb{x})=k.\]
\sssc{Contour map or contour plot}
Contour map or contour plot is the graphical representation of the level curves of a function of two variable, typical represented with a collection of level curves labelled with corresponding levels that are often of same difference relative to their adjacent curves, or with color or height ($z$-axis) showing function value at all points.
\sssc{Level surface}
For a real-valued function $f$ of three variables, the level surface of $f$ at level $k\in\bbR$ is defined as the set of $\mb{x}$ such that
\[f(\mb{x})=k.\]
\sssc{Gradient perpendicular to tangent of level}
For function $F(\mb{x})\colon D\subseteq\bbR^n\to\bbR$, suppose $\mb{r}(t)$ is any parameterized curve lies on $F(\mb{x})=k$, then
\[\nabla F\cdot\mb{r}'=0.\]
\sssc{Tangent flat}
For function $F(\mb{x})\colon D\subseteq\bbR^n\to\bbR$, for each point $\mb{x}_0\in D$, the tangent flat of $F$ at it is given by
\[\nabla F\cdot(\mb{x}-\mb{x}_0)=0.\]
\sssc{Normal line}
For function $F(\mb{x})\colon D\subseteq\bbR^n\to\bbR$, for each point $\mb{x}_0\in D$, the normal line of $F$ at it is given by
\[\mb{x}(t)=t\nabla F(\mb{x}_0)+\mb{x}_0,\quad t\in\bbR.\]
\subsection{Line integral (線積分) or Path integral (路徑積分)}
\subsubsection{Definition}
For a parameterized simple curve $\gamma(t)\colon[a,b]\to\bbR^n$ that is piecewise $C^1$, the line integral of a function $f$ defined on the image of $\gamma$ over $\gamma$ is defined as the Riemann–Stieltjes integral:
\[\int_{\gamma}f(\mb{x})\dd{\mb(x)}\coloneq\int_{t=a}^bf(\gamma(t))\dd{\gamma(t)},\]
In which $\gamma$ is called the integral path.

The integral does not depend on the parametrization of the curve as long as the curve starts at $\gamma(a)$ and ends at $\gamma(b)$ and is piecewise $C^1$.
\subsubsection{Conservative field (保守場)}
A field $f$ whose domain is a subset $U$ of a Euclidean tensor space is called a conservative field iff for all paths $C$ between point $a$ and $b$, the integral 
\[\int_Cf(\mathbf{x})\cdot\mathrm{d}\mathbf{x}\]
are the same, iff for any closed path $C$, 
\[\oint_Cf(\mathbf{x})\cdot\mathrm{d}\mathbf{x}=0.\]
\bit
\item Field $f$ whose domain is $U\subseteq\bbR$ is conservative iff it is continuous.
\item If field $f=(f_x,f_y)$ whose domain is $U\subseteq\bbR^2$ is conservative, then
\[\pdv{f_y}{x}-\pdv{f_x}{y}=0\]
wherever it is defined.
\item Field $f=(f_x,f_y)$ whose domain is a simply connected set $U\subseteq\bbR^2$ is conservative iff
\[\pdv{f_y}{x}-\pdv{f_x}{y}=0\]
for all $(x,y)\in U$.
\item If field $f$ whose domain is $U\subseteq\bbR^3$ is conservative, then
\[\nabla\times f=0\]
wherever it is defined.
\item Field $f$ whose domain is a simply connected set $U\subseteq\bbR^3$ is conservative iff
\[\nabla\times f(\mb{x})=0\]
for all $\mb{x}\in U$.
\eit
\ssc{Multiple integral}
\sssc{Lebesgue integral}
Defined via product measure space.
\sssc{Riemann integral}
For rectangle
\[R=\prod_{i=1}^n[x_{i0},x_{i1}],\]
and function $f$ bounded on $R$,
\[\int_Rf(\mb{x})\dd{V}\coloneq\lim_{(m_1,m_2,\ldots,m_n)\to\infty}\qty(\sum_{j_i=1}^{m_i})_{i=1}^nf\qty(\mb{x}_{(j_i)_{i=1}^n}*)\Delta V,\]
called multiple Riemann sum, where
\[\Delta x_i=\frac{\qty(x_{i1}-x_{i0})}{m_i},\]
\[\Delta V=\prod_{i=1}^n\Delta x_i,\]
and $\mb{x}_{(j_i)_{i=1}^n}*$ is an arbitrary chosen sample point in
\[\prod_{i=1}^n\qty[x_{i0}+\qty(j_i-1)\Delta x_i,x_{i0}+j_i\Delta x_i.\]
For arbitrary subset $E\subseteq R$ and function $g$ bounded on $E$, let $h$ be the trivial extension of $g_E$ to $R$,
\[\int_Eg(\mb{x})\dd{V}\coloneq\int_Rh(\mb{x})\dd{V}.\]
\sssc{Iterated integral}
An integral whose integrand is an integral.
\sssc{Type I and Type II region in two dimensions}
A Type I region is a subset of $\bbR^2$ of the form
\[D=\{(x,y)\mid a\le x\le b\land g_1(x)\le y\le g_2(x)\}\]
where $g_1,g_2$ are continuous functions defined on $[a,b]$.

A Type II region is a subset of $\bbR^2$ of the form
\[D=\{(x,y)\mid a\le y\le b\land g_1(y)\le x\le g_2(y)\}\]
where $g_1,g_2$ are continuous functions defined on $[a,b]$.